\documentclass[10pt]{article}
\usepackage[T1]{fontenc}
\usepackage[utf8]{inputenc}
\usepackage[french]{babel}
\usepackage{amsmath,amssymb,mathrsfs,mathpazo}
\usepackage{geometry}
\geometry{a4paper,margin=2cm}
\usepackage{graphicx,array,tabularx,tabularray,colortbl,multirow,multicol,booktabs,enumitem,pifont,ulem,fancybox,diagbox,cancel}
\usepackage{pstricks,pstricks-add,pst-plot,pst-tree,pst-node,pst-text,pst-func,pst-math,pst-bspline,pst-3dplot,pst-eucl,pst-solides3d}
\let\pscalculate\undefined
\usepackage{pst-bezier}
\usepackage{tikz,pgfplots,tkz-tab}
\usetikzlibrary{arrows,arrows.meta,patterns,calc,positioning,trees,decorations.markings,shapes.geometric,shapes.misc,shapes.arrows,intersections,decorations.pathmorphing,decorations.pathreplacing,backgrounds,shadows,babel}
\pgfplotsset{compat=1.18}
\usepgfplotslibrary{fillbetween}
\usepackage[np]{numprint}
\usepackage{listings,makecell,pas-tableur,twoopt}
\let\barre\undefined
\usepackage[pstricks]{tablvar}
\definecolor{darkgreen}{RGB}{0,100,0}
\providecommand{\red}{\color{red}}
\providecommand{\blue}{\color{blue}}
\providecommand{\green}{\color{green}}
\providecommand{\black}{\color{black}}
\providecommand{\white}{\color{white}}
\providecommand{\gray}{\color{gray}}
\providecommand{\cyan}{\color{cyan}}
\providecommand{\magenta}{\color{magenta}}
\providecommand{\yellow}{\color{yellow}}
\providecommand{\violet}{\color{violet}}
\providecommand{\orange}{\color{orange}}
\providecommand{\euro}{\texteuro}
\providecommand{\EUR}{\texteuro}
\providecommand{\ieme}{\textsuperscript{e}}
\providecommand{\degres}{\ensuremath{^\circ}}
\providecommand{\degree}{\ensuremath{^\circ}}
\providecommand{\vv}[1]{\overrightarrow{#1}}
\providecommand{\Coord}[1]{#1}
\providecommand{\decosix}{}
\providecommand{\starredbullet}{\ensuremath{\bullet}}
\pagestyle{empty}
\setlength{\parindent}{0pt}
\usepackage{adjustbox}\DeclareUnicodeCharacter{2236}{:}\DeclareUnicodeCharacter{2190}{\ensuremath{\leftarrow}}
\providecommand{\R}{\mathbb{R}}
\providecommand{\N}{\mathbb{N}}
\providecommand{\Z}{\mathbb{Z}}
\providecommand{\Q}{\mathbb{Q}}
\providecommand{\D}{\mathbb{D}}
\providecommand{\C}{\mathbb{C}}
\providecommand{\e}{{\mathrm{e}}}
\providecommand{\vect}[1]{\overrightarrow{#1}}
\providecommand{\vectt}[1]{\overrightarrow{\mathrm{#1}}}
\providecommand{\vvt}[1]{\overrightarrow{\mathrm{#1}}}
\providecommand{\np}[2][]{{#2~#1}}
\providecommand{\vv}[1]{{\overrightarrow{#1}}}
\renewcommand{\d}{{\mathrm{d}}}
\providecommand{\degres}{{^\circ}}
\providecommand{\degre}{{^\circ}}
\providecommand{\blue}{\color{blue}}
\providecommand{\red}{\color{red}}
\providecommand{\enskip}{\quad}
\providecommand{\llbracket}{\lbrack\!\lbrack}
\providecommand{\rrbracket}{\rbrack\!\rbrack}
\providecommand{\euro}{{\text{€}}}
\providecommand{\no}{{\text{n°}}}
\providecommand{\interval}[2]{{[#1;#2]}}
\providecommand{\intervFO}[2]{{[#1;#2[}}
\providecommand{\dotfill}{\ldots}

\providecommand{\ds}{\displaystyle}
\providecommand{\textding}[1]{#1}
\providecommand{\Oij}{\(\mathrm{O};\vec{\imath},\vec{\jmath}\)}
\providecommand{\Oijk}{\(\mathrm{O};\vec{\imath},\vec{\jmath},\vec{k}\)}
\providecommand{\Ouv}{\(\mathrm{O};\vec{u},\vec{v}\)}
\providecommand{\intervFF}[2]{\ensuremath{[#1;#2]}}
\renewcommand{\euro}{\ifmmode\text{\texteuro}\else\texteuro\fi}
\DeclareUnicodeCharacter{00B0}{\ifmmode^{\circ}\else\textdegree\fi}
\DeclareUnicodeCharacter{20AC}{\ifmmode\text{\texteuro}\else\texteuro\fi}
% Énoncés : Éducation nationale. Transcription et corrigés : APMEP, Denis Vergès et contributeurs. Découpage et classement : CoopMaths, Sébastien Lozano, Éric Elter et contributeurs. Adaptation de la présentation : Maths Explorer.
% CC BY-SA : https://coopmaths.fr/www/about/
% https://coopmaths.fr/alea/static/e3c/2020/tex/e3c_2020_00_sujet45_2_cor.tex
\begin{document}
\noindent\textit{APMEP / CoopMaths. CC BY-SA. Présentation adaptée par Maths Explorer.}\par\medskip
\subsection*{1.}

La probabilité de tirer une boule rouge est :
\[
p(R) = \dfrac{2}{2+3} = \dfrac{2}{5} = \dfrac{4}{10} = 0{,}4.
\]
La probabilité de tirer une boule noire est :
\[
p(N) = \dfrac{3}{2+3} = \dfrac{3}{5} = \dfrac{6}{10} = 0{,}6.
\]

On peut donc dresser l’arbre pondéré :
\begin{center}
\begin{tikzpicture}[thick, scale=1.5]
\node (P_-1_0) at (-2,-1.5) {$\phantom{A}$};
\node (P_0_0) at (0,-0.5) {$R$};
\draw (P_-1_0) -- (P_0_0) node[midway, above] {$\frac25$};
\node (P_1_0) at (2,-0) {$R$};
\draw (P_0_0) -- (P_1_0) node[midway, above] {$\frac25$};
\node (P_1_1) at (2,-1) {$N$};
\draw (P_0_0) -- (P_1_1) node[midway, below] {$\frac35$};
\node (P_0_2) at (0,-2.5) {$N$};
\draw (P_-1_0) -- (P_0_2) node[midway, below] {$\frac35$};
\node (P_1_2) at (2,-2) {$R$};
\draw (P_0_2) -- (P_1_2) node[midway, above] {$\frac25$};
\node (P_1_3) at (2,-3) {$N$};
\draw (P_0_2) -- (P_1_3) node[midway, below] {$\frac35$};
\end{tikzpicture}
\end{center}

\subsection*{2.}

On a :
\[
p(R \cap R) = p(R) \times p_R(R) = \dfrac{2}{5} \times \dfrac{2}{5} = \dfrac{4}{25} = 0{,}16.
\]

\subsection*{3.}

La probabilité de tirer deux boules noires est :
\[
p(N \cap N) = p(N) \times p_N(N) = \dfrac{3}{5} \times \dfrac{3}{5} = \dfrac{9}{25} = \dfrac{36}{100} = 0{,}36.
\]
Dans ce cas \(X = -20\).

La probabilité de tirer une boule de chaque couleur est donc :
\begin{align*}
1 - p(R \cap R) - p(N \cap N) &= 1 - 0{,}16 - 0{,}36 \\
&= 1 - 0{,}52 \\
&= 0{,}48.
\end{align*}
Dans ce cas \(X = 40\), le dernier cas étant \(X = 20 - 10 = 10\).

D'où le tableau de la loi de probabilité de \(X\) :
\[
\begin{array}{|c|c|c|c|}
\hline
X = x_i& 40 & -20 & 10 \\
\hline
p(X = x_i) & 0{,}16 & 0{,}36 & 0{,}48 \\
\hline
\end{array}
\]

\subsection*{4.}

La probabilité de gagner de l’argent (\(X > 0\)) est égale à \(0{,}16 + 0{,}48 = 0{,}64\).

\subsection*{5.}

L’espérance mathématique de la variable aléatoire \(X\) est :
\begin{align*}
E(X) &= 40 \times 0{,}16 - 20 \times 0{,}36 + 10 \times 0{,}48 \\
&= 6{,}4 - 7{,}2 + 4{,}8 \\
&= 11{,}2 - 7{,}2 \\
&= 4.
\end{align*}

En moyenne, sur un grand nombre de parties, le gain est de 4 € par partie.

\end{document}
